\documentclass[12pt,a4paper]{article}
\usepackage[utf8]{inputenc}
\usepackage{amsmath, amssymb}
\usepackage{enumitem}
\usepackage[margin=2cm]{geometry}

\begin{document}

\begin{center}
    \textbf{\Large Latihan Soal Basa Sunda SMK Kelas X} \\
    \textbf{Persiapan Ujian Akhir Semester (UAS)}
\end{center}
\vspace{0.5cm}

\begin{enumerate}[leftmargin=*]

% === BAGIAN 1: METODOLOGI WAWANCARA ===
\item Tanya jawab langsung pikeun nyajaring informasi ti narasumber disebut na \dots
    \begin{enumerate}[label=\alph*.]
        \item Wawancara
        \item Biantara
        \item Sawala
        \item Dongeng
        \item Sinden
    \end{enumerate}

\item Kagiatan wawancara anu pananya teu dirumuskeun sacara tinulis sateuacanna disebutna wawancara \dots
    \begin{enumerate}[label=\alph*.]
        \item Terstruktur
        \item Bebas
        \item Katutup
        \item Kelompok
        \item Dinas
    \end{enumerate}

\item Jalma anu bertindak salaku méré katerangan, data, fakta, atawa opini dina wawancara disebut \dots
    \begin{enumerate}[label=\alph*.]
        \item Pewawancara
        \item Moderator
        \item Narasumber
        \item Jurnalis
        \item Panyajak
    \end{enumerate}

\item Jalma anu miboga pancén pikeun ngadali prosés tanya jawab sarta ngajukeun patanyaan disebut \dots
    \begin{enumerate}[label=\alph*.]
        \item Pewawancara
        \item Informan
        \item Panata harta
        \item Girang serat
        \item Panumbu catur
    \end{enumerate}

\item Wawancara anu daptar pertanyaanana geus disusun sacara rapih jeung sistematis saacan panggih jeung narasumber disebut \dots
    \begin{enumerate}[label=\alph*.]
        \item Wawancara bebas
        \item Wawancara mutlak
        \item Wawancara dinas
        \item Wawancara terstruktur
        \item Wawancara lisan
    \end{enumerate}

\item Kecap pananya dina basa Sunda anu dipaké pikeun nanyakeun tokoh, jalma, atawa palaku nyaéta \dots
    \begin{enumerate}[label=\alph*.]
        \item Naon
        \item Saha
        \item Iraha
        \item Naha
        \item Mana
    \end{enumerate}

\item Upami urang hoyong nanyakeun alesan atawa kasang tukang (latar belakang) hiji kajadian, kecap pananya anu merenah nyaéta \dots
    \begin{enumerate}[label=\alph*.]
        \item Kumaha
        \item Sabaraha
        \item Naha
        \item Mana
        \item Naon
    \end{enumerate}

\item Ragam basa anu wajib dipaké nalika ngawawancara narasumber anu leuwih kolot atawa tokoh nu dipihormat nyaéta \dots
    \begin{enumerate}[label=\alph*.]
        \item Basa Loma
        \item Basa Kasar
        \item Basa Hormat (Lemes)
        \item Basa Cohag
        \item Basa Garihal
    \end{enumerate}

\item Basa loma atawa akrab dina kagiatan wawancara ukur bisa dipaké nalika narasumberna \dots
    \begin{enumerate}[label=\alph*.]
        \item Kepala Sakola
        \item Guru produktif TKJ
        \item Babaturan saumuran nu geus akrab
        \item Tokoh masyarakat
        \item Pejabat pamaréntah
    \end{enumerate}

\item Tujuan utama saurang siswa ngalakukeun kagiatan wawancara nyaéta \dots
    \begin{enumerate}[label=\alph*.]
        \item Neangan hiburan
        \item Nyajaring informasi sarta data valid
        \item Pamer kaahlian ngomong
        \item Nyieun carita fiksi
        \item Ngadebat narasumber
    \end{enumerate}

\item Nangtukeun téma, milih narasumber, sarta ngontak narasumber pikeun nyieun jangji kaasup kana tahapan \dots
    \begin{enumerate}[label=\alph*.]
        \item Tahap palaksanaan
        \item Tahap nyusun laporan
        \item Tahap tatahar (persiapan)
        \item Tahap evaluasi akhir
        \item Tahap publikasi
    \end{enumerate}

\item Kecap pananya 'Iraha' dipaké ku pewawancara pikeun ménta katerangan ngeunaan \dots
    \begin{enumerate}[label=\alph*.]
        \item Tempat kajadian
        \item Waktu lumangsungna peristiwa
        \item Jumlah barang
        \item Cara ngaréngsékeun masalah
        \item Identitas narasumber
    \end{enumerate}

\item Wawancara anu ngasupkeun gurat badag (outline) topik tapi patanyaanana dimekarkeun sacara bébas di lapangan disebut \dots
    \begin{enumerate}[label=\alph*.]
        \item Wawancara bebas terpimpin
        \item Wawancara mutlak terikat
        \item Wawancara rahasia
        \item Wawancara kelompok
        \item Wawancara spontan
    \end{enumerate}

\item Prosés ngolah data tina rékaman atawa catetan sarta nuliskeunana deui dina wangun artikel kaasup kana \dots
    \begin{enumerate}[label=\alph*.]
        \item Tahap tatahar
        \item Tahap pelaksanaan
        \item Tahap nyusun laporan
        \item Tahap konfirmasi awal
        \item Tahap pemantauan
    \end{enumerate}

\item Kecap 'naroskeun' mangrupa wangun ragam basa lemes tina kecap \dots
    \begin{enumerate}[label=\alph*.]
        \item Nyarita
        \item Menta
        \item Nanya
        \item Ngadenge
        \item Nempo
    \end{enumerate}

% === BAGIAN 2: SASTRA BUHUN CARITA BABAD ===
\item Carita atawa dongeng anu ngandung unsur sajarah lokal disebut \dots
    \begin{enumerate}[label=\alph*.]
        \item Mitos
        \item Fabel
        \item Babad
        \item Roman
        \item Parabel
    \end{enumerate}

\item Carita babad teh kaasup kana wangun carita buhun anu eusina ngenaan \dots
    \begin{enumerate}[label=\alph*.]
        \item Kahirupan teknologi modern
        \item Kajadian atawa lalakon di jaman baheula
        \item Dongeng sato anu bisa nyarita
        \item Kritik sosial jaman ayeuna
        \item Siloka kaagamaan modern
    \end{enumerate}

\item Nu jadi ciri khas tina carita babad nyaeta \dots
    \begin{enumerate}[label=\alph*.]
        \item Tokohna sato kabeh
        \item Ditulis dina wangun novel modern
        \item Miboga unsur sajarah lokal
        \item Teu boga amanat carita
        \item Pengarangna kapanggih jaman ayeuna
    \end{enumerate}

\item Utama tina dijieunna atawa disebarkeunana carita babad, sabenerna tujuan carita babad nyaeta \dots
    \begin{enumerate}[label=\alph*.]
        \item Jang néangan kauntungan duit
        \item Nyalamat sajarah baheula
        \item Ngabobodo masyarakat leutik
        \item Nyebarkeun warta bohong (hoax)
        \item Ngarobah sajarah asli karaton
    \end{enumerate}

\item Dina carita babad, biasana aya unsur nu dicampur antara sajarah jeung \dots
    \begin{enumerate}[label=\alph*.]
        \item Pamohalan
        \item Teknologi
        \item Statistik
        \item Jurnalistik
        \item Berita ilmiah
    \end{enumerate}

\item Anu kaasup kana tilu tihang utama ciri ciri carita babad nyaéta \dots
    \begin{enumerate}[label=\alph*.]
        \item Fabel, parabel, legenda
        \item Novel, sajak, biantara
        \item Lalakon baheula, pamohalan, miboga unsur sajarah
        \item Tokoh sato, nyata, modern
        \item Sajarah asli, ilmiah, faktual kabeh
    \end{enumerate}

\item Runtuyan pirang pirang kajadian nu ngawangun dina carita babad disebut \dots
    \begin{enumerate}[label=\alph*.]
        \item Téma
        \item Palaku
        \item Latar
        \item Galur
        \item Amanat
    \end{enumerate}

\item Karajaan munggaran di jawa barat dumasar catetan sajarah lokal jeung babad nyaeta \dots
    \begin{enumerate}[label=\alph*.]
        \item Pajajaran
        \item Tarumanagara
        \item Galuh
        \item Sumedang Larang
        \item Cirebon
    \end{enumerate}

\item Carita babad anu nyaritakeun ngeunaan runtagna Karajaan Pajajaran nyaéta \dots
    \begin{enumerate}[label=\alph*.]
        \item Babad Limbangan
        \item Babad Cirebon
        \item Babad Bogor
        \item Babad Panjalu
        \item Babad Banten
    \end{enumerate}

\item Istilah 'pamohalan' dina carita babad miboga harti \dots
    \begin{enumerate}[label=\alph*.]
        \item Hal anu fiksi ilmiah
        \item Hal-hal anu teu asup akal atawa magis
        \item Catetan sajarah sagemblengna bener
        \item Tokoh anu bener-bener aya jaman ayeuna
        \item Tempat anu can kungsi kapanggih
    \end{enumerate}

\item Tokoh raja anu kacida kawentarna ti Karajaan Tarumanagara dumasar prasasti jeung babad nyaéta \dots
    \begin{enumerate}[label=\alph*.]
        \item Prabu Siliwangi
        \item Sunan Gunung Jati
        \item Purnawarman
        \item Prabu Borosngora
        \item Jayadewata
    \end{enumerate}

\item Galur anu nyaritakeun kajadian dumasar urutan waktu ti mimiti nepi ka tungtung disebut \dots
    \begin{enumerate}[label=\alph*.]
        \item Galur ti balik (Mundur)
        \item Galur merelé (Maju)
        \item Galur campuran
        \item Galur rumpang
        \item Galur bebas
    \end{enumerate}

\item Babad Galuh nyaéta carita babad anu nyaritakeun sajarah lokal di wewengkon \dots
    \begin{enumerate}[label=\alph*.]
        \item Garut
        \item Ciamis
        \item Sumedang
        \item Bogor
        \item Banten
    \end{enumerate}

\item Unsur intrinsik babad anu mangrupa gagasan poko atawa ide pusat anu ngajiwaan sakabéh eusi carita disebut \dots
    \begin{enumerate}[label=\alph*.]
        \item Latar
        \item Palaku
        \item Téma
        \item Galur
        \item Mimik
    \end{enumerate}

\item Tempat, waktu, jeung suasana anu dicaritakeun dina carita babad kaasup kana unsur \dots
    \begin{enumerate}[label=\alph*.]
        \item Téma
        \item Latar (Setting)
        \item Karakter
        \item Gaya basa
        \item Sudut pandang
    \end{enumerate}

% === BAGIAN 3: SISTEM KEKERABATAN PANCAKAKI ===
\item Naon anu dimaksud pancakaki teh \dots
    \begin{enumerate}[label=\alph*.]
        \item Aturan ngigel Sunda
        \item Silsilah kaluwarga
        \item Cara masak kadaharan tradisional
        \item Waditra kasenian degung
        \item Elmu beladiri tatar Sunda
    \end{enumerate}

\item Dina rundayan pancakaki ka luhur, kolot na buyut sok disebut \dots
    \begin{enumerate}[label=\alph*.]
        \item Aki
        \item Janggawareng
        \item Bao
        \item Udeg-udeg
        \item Gantung siwur
    \end{enumerate}

\item Istilah bau bau sinduk dina pancakaki hartina \dots
    \begin{enumerate}[label=\alph*.]
        \item Baraya keneh sanajan tos laen
        \item Dulur anu saindung sabapa
        \item Anakna ti lanceuk urang
        \item Kolotna ti aki jeung nini urang
        \item Baraya anu geus tilar dunya
    \end{enumerate}

\item Istilah dulur pet kuhinis dina pancakaki hartina \dots
    \begin{enumerate}[label=\alph*.]
        \item Dulur angkat
        \item Dulur sabapa wungkul
        \item Dulur saindung wungkul
        \item Dulur kandung (saindung sabapa)
        \item Dulur sabaraya jauh
    \end{enumerate}

\item Gantung siwur dina pancakaki nyaeta bapana \dots
    \begin{enumerate}[label=\alph*.]
        \item Buyut
        \item Bao
        \item Janggawareng
        \item Udeg udeg
        \item Anak
    \end{enumerate}

\item Ka anak lanceuk disebut alo, ari ka anak adi disebutna \dots
    \begin{enumerate}[label=\alph*.]
        \item Suwar
        \item Kapiadi
        \item Kapilanceuk
        \item Incu
        \item Bao
    \end{enumerate}

\item Jalur silsilah kulawarga nangtung anu nataan katurunan saluhureun kolot urang disebut rundayan \dots
    \begin{enumerate}[label=\alph*.]
        \item Ka handap
        \item Ka gigir
        \item Ka luhur (Ascending)
        \item Sampingan
        \item Kahuripan
    \end{enumerate}

\item Katurunan langsung kahiji ti urang sarta salaki/pamajikan disebut \dots
    \begin{enumerate}[label=\alph*.]
        \item Incu
        \item Anak
        \item Umpi
        \item Cicip
        \item Bao
    \end{enumerate}

\item Kolotna ti bapa atawa indung urang dina tatanan pancakaki disebut \dots
    \begin{enumerate}[label=\alph*.]
        \item Buyut
        \item Bao
        \item Aki / Nini
        \item Janggawareng
        \item Gantung siwur
    \end{enumerate}

\item Anakna ti incu urang dina rundayan pancakaki ka handap disebut \dots
    \begin{enumerate}[label=\alph*.]
        \item Umpi
        \item Cicip
        \item Munyang
        \item Gewari
        \item Saompi
    \end{enumerate}

\item Kolotna ti bao urang (tingkat kalima ka luhur) disebut \dots
    \begin{enumerate}[label=\alph*.]
        \item Buyut
        \item Janggawareng
        \item Udeg-udeg
        \item Gantung siwur
        \item Kolot
    \end{enumerate}

\item Tatanan pancakaki ka luhur saatos Buyut nyaéta \dots
    \begin{enumerate}[label=\alph*.]
        \item Aki
        \item Udeg-udeg
        \item Bao
        \item Gantung siwur
        \item Umpi
    \end{enumerate}

\item Dulur anu kaluar tina rahim indung sarta bapa anu sarua disebut ogé \dots
    \begin{enumerate}[label=\alph*.]
        \item Bau-bau sinduk
        \item Dulur pet ku hinis
        \item Kapilanceuk
        \item Kapiadi
        \item Alo
    \end{enumerate}

\item Dina rundayan ka handap, anakna ti umpi urang disebut \dots
    \begin{enumerate}[label=\alph*.]
        \item Incu
        \item Cicip
        \item Munyang
        \item Gewari
        \item Saompi
    \end{enumerate}

\item Naon fungsi sosial utama urang nyaho kana pancakaki kulawarga \dots
    \begin{enumerate}[label=\alph*.]
        \item Sangkan bisa pamer katurunan
        \item Sangkan teu pareumeun obor (kaleungitan baraya)
        \item Jang néangan warisan dukun
        \item Sangkan kasohor di media sosial
        \item Pikeun nangtukeun kabeungharan kulawarga
    \end{enumerate}

% === BAGIAN 4: ANAlISIS PUISI MODERN (SAJAK) ===
\item Naon anu dimaksud sajak teh \dots
    \begin{enumerate}[label=\alph*.]
        \item Dongeng buhun
        \item Prosa luyur
        \item Puisi wangun bebas
        \item Artikel berita
        \item Sisindiran
    \end{enumerate}

\item Nu disebut amanan dina sajak teh hartina \dots
    \begin{enumerate}[label=\alph*.]
        \item Pesan moral / talatah nasehat
        \item Jumlah kecap dina sabaris
        \item Musikalisasi sora
        \item Ngaran panyajak asli
        \item Pilihan kecap konotatif
    \end{enumerate}

\item Gaya bahasa anu sok dipake dina sajak biasana ngagunakeun \dots
    \begin{enumerate}[label=\alph*.]
        \item Bahasa ilmiah baku
        \item Bahasa kasar pisan
        \item Bahasa endah (estetis/kiasan)
        \item Bahasa deungeun (Asing)
        \item Bahasa pemrograman komputer
    \end{enumerate}

\item Sebutkeun salasahiji tujuan utama dina nyieun sajak nyaéta \dots
    \begin{enumerate}[label=\alph*.]
        \item Jang ngeusian kolom iklan berita
        \item Jang nyatakeun rasa, pikiran jeung pangalaman batin
        \item Pikeun ngabobodo nu maca karyana
        \item Jang nyieun téks biantara resmi
        \item Pikeun ngitung guru wilangan pupuh
    \end{enumerate}

\item Naon ari anu disebut padalisan teh \dots
    \begin{enumerate}[label=\alph*.]
        \item Musik sajak
        \item Judul sajak
        \item Bait sajak
        \item Diksi kiasan
        \item Rima tungtung
    \end{enumerate}

\item Saha ngaran panyajak anu kasohor, jeung dianggap salasahiji seniman sajak sastra sunda nu nyebarkeun sajak dina mangsa awal \dots
    \begin{enumerate}[label=\alph*.]
        \item Ajip Rosidi
        \item Kis WS
        \item Yus Rusyana
        \item Wahyu Wibisana
        \item Rafles Lesmana
    \end{enumerate}

\item Samemeh rek nulis sajak aya sabaraha hal anu kudu di cangkumm (dikuasai) \dots
    \begin{enumerate}[label=\alph*.]
        \item 3 hal
        \item 4 hal
        \item 5 hal
        \item 6 hal
        \item 7 hal
    \end{enumerate}

\item Naon ari anu disebut imaji teh \dots
    \begin{enumerate}[label=\alph*.]
        \item Gambar poto panyajak
        \item Citraan atawa gambaran angen-angen/pikiran
        \item Aturan sora tungtung kalimah
        \item Buku kumpulan puisi modern
        \item Téknik maca bedas di panggung
    \end{enumerate}

\item Sajak mangrupa wangun puisi anu sipatna bébas, hartina nyaéta \dots
    \begin{enumerate}[label=\alph*.]
        \item Bébas nulis kasar sacara jorang
        \item Henteu kabeungkeut ku aturan pupuh saperti guru lagu jeung wilangan
        \item Teu kudu boga téma jeung amanat
        \item Bisa ditulis ku saha waé tanpa ngaran asli
        \item Bébas dibaca bari lumpat
    \end{enumerate}

\item Sajak mimiti asup jeung mekar dina khazanah sastra Sunda dina sabudeureun taun \dots
    \begin{enumerate}[label=\alph*.]
        \item 1920
        \item 1945
        \item 1946
        \item 1965
        \item 2000
    \end{enumerate}

\item Sikep panyajak kana masalah utama anu aya dina sajakna (contona kaayaan sedih, bungah, keueung) disebut unsur batin \dots
    \begin{enumerate}[label=\alph*.]
        \item Rasa (Feeling)
        \item Nada (Tone)
        \item Diksi
        \item Wirahma
        \item Padalisan
    \end{enumerate}

\item Sikep panyajak ka jalma anu maca karyana (saperti nada mapatahan, nyindir, atawa ngajak) disebut \dots
    \begin{enumerate}[label=\alph*.]
        \item Téma
        \item Nada (Tone)
        \item Rasa
        \item Imaji
        \item Kalimat
    \end{enumerate}

\item Unsur panyusun sajak anu bisa ngahudangkeun panca indra titingalan (visual) sangkan nu maca saolah nempo kajadian disebut \dots
    \begin{enumerate}[label=\alph*.]
        \item Imaji titingalan (visual)
        \item Imaji déngéan (auditif)
        \item Imaji rampaan (taktil)
        \item Diksi konotatif
        \item Amanat transparan
    \end{enumerate}

\item Ngaran lengkep ti tokoh panyajak Sunda anu kawentar \textit{Kis WS} nyaéta \dots
    \begin{enumerate}[label=\alph*.]
        \item Kiswanda Sastradinata
        \item Kisyandi Wira Sudarma
        \item Kiswandi Sastra Wijaya
        \item Kiswan Supriatna
        \item Kismandoko Sastro
    \end{enumerate}

\item Perkara pangheulana anu kudu dicangking dumasar 6 léngkah nulis sajak nyaéta \dots
    \begin{enumerate}[label=\alph*.]
        \item Milih diksi basa éndah
        \item Ngatur wirahma sora
        \item Nangtukeun téma (ide dasar sajak)
        \item Nyusun amanat talatah
        \item Ngolah imaji citraan
    \end{enumerate}

\end{enumerate}
\end{document}